% 主文档配置：电路理论基础笔记
\documentclass[12pt,a4paper,openany]{book}  % openany取消章节强制从奇数页开始
% ===================== 核心包引入 =====================
% 中文支持（必须用XeLaTeX编译）
\usepackage[UTF8]{ctex}
% 数学公式与符号支持
\usepackage{amsmath,amssymb}
% 颜色支持（封面/高亮使用）
\usepackage{xcolor}
% 页面布局设置
\usepackage{geometry}
% 标题格式自定义
\usepackage{titlesec}
% 分页/环境控制工具
\usepackage{etoolbox}
% 目录样式优化
\usepackage{tocloft}
% 页码样式支持
\usepackage{fancyhdr}
% 关键：添加超链接/目录跳转支持
\usepackage{hyperref}       % 核心跳转包
\usepackage{bookmark}       % 优化PDF书签（可选但推荐）

% ===================== 超链接配置（解决跳转问题） =====================
% 超链接样式（可选，可自定义颜色/样式）
\hypersetup{
    colorlinks=true,        % 链接文字上色（false为边框）
    linkcolor=blue,         % 内部链接（目录→正文）颜色
    anchorcolor=black,      % 锚点颜色
    citecolor=green,        % 引用颜色
    urlcolor=cyan,          % 网址颜色
    pdfstartview=FitH,      % PDF打开时适配宽度
    % bookmarks=true,       % 已被bookmark包自动启用，删除此行
    bookmarksnumbered=true, % 书签显示章节编号
    pdfpagemode=UseOutlines % 打开PDF时显示书签栏
}



% ===================== 全局配置 =====================
% 页面边距（A4纸，适配阅读）
\geometry{
    left=2.5cm,
    right=2.5cm,
    top=3cm,
    bottom=3cm,
    headsep=10pt,  % 页眉间距
    footskip=15pt  % 页脚间距
}

% 标题格式优化（适配中文排版）
\titleformat{\chapter}
    {\centering\Large\bfseries\color{black}}  % 居中、大号、加粗、黑色
    {第\thechapter\ 章}                      % 章节前缀
    {1em}                                     % 前缀与标题间距
    {}                                        % 标题内容格式

\titleformat{\section}
    {\large\bfseries\color{black}}
    {\thesection}
    {1em}
    {}

\titleformat{\subsection}
    {\bfseries\color{black}}
    {\thesubsection}
    {1em}
    {}

% 标题间距调整（消除多余空白）
\titlespacing{\chapter}{0pt}{-20pt}{20pt}  % 章节标题上下间距
\titlespacing{\section}{0pt}{12pt}{6pt}    % 小节标题上下间距
\titlespacing{\subsection}{0pt}{8pt}{4pt}   % 子小节标题上下间距

% ===================== 分页/目录修复 =====================
% 取消目录自动分页，消除封面后空白页
\pretocmd{\tableofcontents}{\clearpage\thispagestyle{empty}}{}{}

% 修复LaTeX默认的分页写入逻辑，避免空白页
\makeatletter
\let\originalclearpage\clearpage
\renewcommand{\clearpage}{%
    \originalclearpage
    \if@filesw
        \immediate\write\@auxout{\string\@setckpt{\@currenvir}{\thepage}}%
    \fi
}
\makeatother

% ===================== 公式注释快捷命令 =====================
% 统一公式中文注释样式（与子文件兼容）
\newcommand{\formulaNote}[1]{\quad \text{#1}}

% ===================== 文档主体 =====================
\begin{document}
% --------------------- 封面与目录 ---------------------
\pagestyle{empty}          % 封面/目录页隐藏页码
% 封面 - 适配book文档类，无编译报错
\begin{titlepage}
    % 全局居中，取消默认缩进
    \centering
    
    % 顶部空行，提升整体布局
    \vspace*{2cm}
    
    % 主标题
    {\Huge \bfseries 电路理论基础笔记} \\
    \vspace{0.8cm}
    % 副标题
    {\Large \textit{核心知识点整理}} \\
    \vspace{0.5cm}
    % 版本标识（依赖main.tex中的xcolor包）
    {\large \textcolor{gray}{0.5残缺版}} \\
    
    % 分割线（提升视觉层次）
    \vspace{2cm}
    \rule{8cm}{0.5pt}
    \vspace{2cm}
    
    % 名言：替换为电路领域知名科学家（尼古拉·特斯拉）
    {\Large \textit{“我知道有一个无限的能量源，它就在我们身边，我们只需要学会如何去捕捉它。”}} \\
    \vspace{1cm}
    {\large —— 尼古拉·特斯拉（Nikola Tesla）}
    
    % 中间大间距，突出底部信息
    \vspace{5cm}
    
    % 作者信息
    {\large 整理人：浮萍} \\
    \vspace{0.5cm}
    % 整理时间
    {\large 整理时间：2026年} \\
    
    % 底部填充，让内容居中偏上
    \vfill
    
    % 页脚小字
    {\small \textcolor{gray}{仅用于学习参考 · 请勿商用}}
\end{titlepage}

% 封面后强制清空分页，避免影响后续排版
\clearpage
\thispagestyle{empty}  % 确保封面后无页码残留              % 引入封面文件
\clearpage                 % 封面后手动分页（仅1次）

\tableofcontents           % 生成目录（现在支持点击跳转）
\thispagestyle{empty}      % 目录页无页码
\clearpage                 % 目录后分页

% --------------------- 正文部分 ---------------------
\pagestyle{plain}          % 正文恢复普通页码样式
\pagenumbering{arabic}     % 正文页码从1开始（阿拉伯数字）

% 引入所有章节子文件（按顺序）
% 第1章 电路元件与电路基本定律（章节子文件，需被主文件调用）
\chapter{电路元件与电路基本定律}

\section{电流、电压与电功率}
\subsection{电流}
荷电质点的有序运动形成电流。为了衡量电流的大小，在时间$dt$内，通过某截面的电荷量的代数和为$dq$，则电流定义为
\[
i=\frac{dq}{dt}
\]
其方向规定为正电荷运动的方向。

电荷的$\mathrm{SI}$(\text{国际单位制})单位为库仑(\text{符号为}$\mathrm{C}$)，时间的$\mathrm{SI}$单位为秒时，电流的$\mathrm{SI}$单位为安培(\text{符号为}$\mathrm{A}$)，$1\mathrm{A}=1\mathrm{C}/1\mathrm{s}$。

电流的方向既可用箭头“$\rightarrow$”表示，也可用双下标表示。例如电流可从a点流向b点或相反，分别用$i_{ab}$和$i_{ba}$表示。在分析复杂电路时，难以判断电流的真实方向，而且有些电路中电流的真实方向随时间不断改变。为此，在列写电路方程时，可任意假设一个电流的方向，称为电流的参考方向，并用箭头标示在图中。选定参考方向后，电流便是代数量。电流的真实方向由电流代数量的符号和参考方向来判断。当按参考方向求得的电流为正值时，表明其真实方向与参考方向一致；否则相反。

量值和方向不随时间变化的电流称为直流。

\subsection{电压}
单位正电荷从电路中的一点移动到另一点，电场力对电荷做功。设电场力做功为$dW_{ab}$，则电压定义为
\[
u_{ab}=\frac{dW_{ab}}{dq}
\]
在式中，$u_{ab}$表示$\text{a、b}$两点间的电压，其物理意义是单位正电荷从$\text{a}$点移动到$\text{b}$点时电场力所做的功。

电压又称电位差，与电压相关的另一个物理量是电位。电位的物理意义是：在电路中可任选一点为参考点(\text{零电位点，常用符号}$\perp$\text{表示})，电路中某点与参考点之间的电压称为该点的电位，用$\phi$表示。

电压的方向规定为从高电位指向低电位的方向。电压参考方向可用双下标表示（如$u_{ab}$表示$\text{a}$点为高电位、$\text{b}$点为低电位）。

若电压的参考方向与真实方向一致，电压为正值；反之则为负值。在分析电路问题时，同样需要先设定电压的参考方向，再通过计算结果的符号判断真实方向。

\subsection{电动势}
从能量的角度看，电源是将其他形式的能量转换为电能的装置。设在$dt$的时间内，一电源使正电荷$dq$从负极经电源内部移至正极所做的功$dW$，电源的电动势定义为
\[
e=\frac{dW}{dq}
\]
即电源电动势的物理意义是单位正电荷从电源负极移至正极时，电源非静电力所做的功。

\subsection{电功率}
电功率是用以衡量电能转换速率的物理量，单位时间内电场力所做的功即为电功率，则电功率定义为
\[
p=\frac{dW}{dt}
\]

当电能被电路吸收时，功率为正值；当电路发出电能时，功率为负值。设某部分电路的电压参考方向与电流参考方向一致（关联参考方向），则该部分电路吸收的功率为$p=ui$；如果所选定的电压和电流的参考方向相反（非关联参考方向），则吸收的功率为$p=-ui$。

功率等于能量的转换速率，根据式可得在$t_0$到$t$时间内，电路吸收（电压、电流为关联参考方向时）或发出的能量为
\[
W(t)=\int_{t_0}^tp(\xi)d\xi=\int_{t_0}^tu(\xi)i(\xi)d\xi
\]
与判断功率正负类似，能量为正值时表示电路吸收能量，负值时表示发出能量。实际工程中，电气设备都有额定功率，如发电站、电动机等，若超过额定值，设备或部件可能因过热损坏。

\section{电阻元件}
在电路模型中，用电阻元件近似表示实际电阻器、灯泡、电炉等耗能器件的电气特性。在电路理论中，将流过相同电流的两个端子称为一个端口。线性二端电阻的端口电压与电流之间遵循欧姆定律，即二者成正比关系。

欧姆定律表达式为
\[
u=Ri
\]
或
\[
i=Gu
\]

若$u\text{、}i$参考方向相反，则欧姆定律应写为$u=-Ri$或$i=-Gu$。电阻与电导是反映同一电阻元件性能而互为倒数的两个参数，即
\[
RG=1
\]
线性电阻的端口特性（$u-i$曲线）是过原点的直线。

如果电阻$R=0$，则通过任意有限电流时其电压恒为零，故相当于短路，即短路线的电阻为零；如果$G=0$（$R\rightarrow\infty$），则施加任意有限电压时其电流恒为零，相当于开路。

电阻吸收的功率是$p=ui$，将欧姆定律代入，得
\[
p=ui=Ri^2=Gu^2
\]

如果$R\text{、}G$是正值，则$p$非负，这样的电阻消耗电功率，属于耗能元件。从能量角度看，直到$t$瞬间输入电阻中的全部电能将是
\[
W=\int_{-∞}^tp(\xi)d\xi=\int_{-∞}^tu(\xi)i(\xi)d\xi=R\int_{-∞}^ti^2(\xi)d\xi=G\int_{-∞}^tu^2(\xi)d\xi
\]
此式就是物理学中的焦耳定律，表明电阻消耗的电能全部转换为热能。

对于非负的电阻，在任一时刻，上式的结果非负。这表明：从全过程来看，这种元件只能吸收电能而不能发出电能。具有此种性质的元件，称为无源元件。

在电路模型中还存在一种负电阻，其$R<0$，能发出电能，属于有源元件（通常由电子电路实现）。

从广义上讲，只要端口电压、电流关系是代数方程，就属于电阻元件，其电压、电流方程一般形式为
\[
u=u(i)
\]
或
\[
i=i(u)
\]
上式中$u(i)$和$i(u)$分别表示电压是电流的函数、电流是电压的函数。如果电压、电流不是正比关系，即电压、电流关系不是过原点的直线，就称为非线性电阻。例如，半导体二极管就属于非线性电阻，其电压、电流关系为$i=I_0(e^{u/U_T}-1)$（$I_0$为反向饱和电流，$U_T$为热电压）。

\section{电容元件}
电容元件的电路符号中，固定电容和可变电容有不同表示。在电路模型中，用电容元件表示电容器的主要电气特性，即电荷与电压的关系。

当在电容器两端通过电源施加电压时，两极板上分别聚集起等量的正、负电荷，并在极板之间建立电场，储存有电场能量。当电源移去后，电荷和电场可以继续存在（理想电容），但实际电容器因漏电流会逐渐放电。电容能够储存电荷和电场能，是储能元件。

从物理学已知，当电容器填充线性介质（即介电常数与电场强度无关，如空气、橡皮、瓷、纱、胶木和云母等）时，它储存的电荷$q$与极板间电压$u$成正比。若电压的参考方向是从正极板指向负极板，则线性电容的$q-u$关系为
\[
q=Cu
\]
式中$C$是电容系数，其单位是$\mathrm{C/V}$(\text{库/伏})，称为法拉(\text{符号为}$\mathrm{F}$)，实际工程中常用微法（$\mu\mathrm{F}$，$1\mu\mathrm{F}=10^{-6}\mathrm{F}$）、皮法（$\mathrm{pF}$，$1\mathrm{pF}=10^{-12}\mathrm{F}$）。

平行板电容器的电容为
\[
C=\frac{\varepsilon A}{d}
\]
式中，$A$为每个极板的有效面积，$d$为极板间距离，$\varepsilon$为介质的介电常数（真空介电常数$\varepsilon_0=8.85\times10^{-12}\mathrm{F/m}$）。对线性介质，电容$C$是常量，称为线性电容；否则称为非线性电容。

线性电容的电荷-电压关系曲线是位于第$\mathrm{I}$、第$\mathrm{III}$象限且经过原点的直线。根据电流的定义，电流等于极板电荷的时间变化率
\[
i=\frac{dq}{dt}
\]
代入$q=Cu$，得
\[
i=C\frac{du}{dt}
\]
即线性电容的端口电流与端口电压的时间变化率成正比。当在端口上施加直流电压时，电压变化率为零，端口电流为零，电容等效成开路。

反过来，电荷$q$等于电流$i$对时间的积分，即
\[
q(t)=\int_{-∞}^ti(\xi)d\xi
\]
式表明$t$时刻电容上的电荷量是此时刻以前由电流充电和放电而“积累”的结果，所以电容具有记忆的性质，属于记忆元件。上式积分下限取$-∞$，意指从很久以前算起，电容上的电荷为零（初始状态）。

由$q=Cu$，得
\[
u(t)=\frac{q(t)}{C}=\frac{1}{C}\int_{-∞}^ti(\xi)d\xi
\]
式即为线性电容的$u-i$关系，式中$u\text{、}i$应为关联参考方向；否则要在式中添一负号（$u(t)=-\frac{1}{C}\int_{-∞}^ti(\xi)d\xi$）。

电阻元件是无记忆元件（电压仅由当前电流决定），而电容是记忆元件（电压与历史电流有关）。

如果研究$t_0$时刻以后的电容电压，可将积分拆分为两部分：
\[
u(t)=\frac{1}{C}\int_{-∞}^{t_0}i(\xi)d\xi+\frac{1}{C}\int_{t_0}^ti(\xi)d\xi=u(t_0)+\frac{1}{C}\int_{t_0}^ti(\xi)d\xi
\]
式中，$u(t_0)$表示$t_0$时刻以前全部电流对电压的影响，称为电容在$t_0$时的初始电压。

输入线性电容的功率为
\[
p=ui=Cu\frac{du}{dt}
\]
根据能量与功率的关系，由式可以求得截止到$t$时刻电容吸收的总能量
\[
W_e(t)=\int_{-∞}^tp(\xi)d\xi=\int_{-∞}^tCu(\xi)\frac{du(\xi)}{d\xi}d\xi=C\int_{u(-∞)}^{u(t)}u du
\]
在$t=-∞$时电容尚未充电，所以$u(-∞)=0$，故电容吸收的总能量为
\[
W_e(t)=\frac{1}{2}Cu^2=\frac{q^2}{2C}
\]
式就是电容储存电场能量的计算公式，表明电容储存的能量仅与当前电压有关，与充电过程无关。

\section{电感元件}
在电机、变压器、电磁铁等电气设备中，线圈是核心部件。尽管实际线圈存在电阻，但其主要电气特性是通过电流产生磁场、储存磁场能量。在电路模型中，用电感元件表示线圈的这一核心特性。

如果线圈的磁场存在于线性介质中（如空气、非铁磁材料），则磁链与电流成正比，称为线性电感；否则（如铁磁材料）称为非线性电感。

设线性电感的磁链$\Psi$与电流$i$的参考方向符合右手螺旋定则，则二者之间的关系为
\[
\Psi=Li
\]
式中$L$为电感系数，磁链的单位是韦伯(\text{符号}$\mathrm{Wb}$)，电流的单位是安培($\mathrm{A}$)时，电感的$\mathrm{SI}$单位是亨利(\text{符号}$\mathrm{H}$)，$1\mathrm{H}=1\mathrm{Wb/A}$，实际常用毫亨（$\mathrm{mH}$）、微亨（$\mu\mathrm{H}$）。

电感的特性还可以用图形表示，与式对应的磁链—电流关系是一条通过$\Psi-i$平面原点的直线且位于第$\mathrm{I}$、第$\mathrm{III}$象限。

根据电磁感应定律，磁链随时间的变化要在电感中产生感应电动势，记作$e$。当电压$u$、电动势$e$与电流$i$取关联参考方向（电压参考方向与电动势相反），则
\[
u=-e=-(-\frac{d\Psi}{dt})=\frac{d\Psi}{dt}
\]
对线性电感，代入$\Psi=Li$，得
\[
u=L\frac{di}{dt}
\]
即线性电感的端口电压与端口电流的时间变化率成正比。当电感中流过直流电流时，电流变化率为零，端口电压为零，电感等效成短路。

如果从电压求磁链或电流，则需积分，即
\[
\Psi(t)=\int_{-∞}^tu(\xi)d\xi=\Psi(t_0)+\int_{t_0}^tu(\xi)d\xi
\]
\[
i(t)=\frac{1}{L}\int_{-∞}^tu(\xi)d\xi=i(t_0)+\frac{1}{L}\int_{t_0}^tu(\xi)d\xi
\]
此两式表明，电感中某一瞬间的磁链和电流决定于此瞬间以前的全部电压，因此电感也属于记忆元件。

上述$u-i$关系要求$u\text{、}i$为关联参考方向；否则要在式中添一负号（$u=-L\frac{di}{dt}$）。

线性电感吸收的功率为
\[
p=ui=Li\frac{di}{dt}
\]
截止到$t$时刻电感吸收的能量为
\[
W_m = \int_{-\infty}^{t} p(\xi) \, d\xi = L \int_{i(-\infty)}^{i(t)} i(\xi) \, di(\xi) = \frac{1}{2} L i^2 \bigg|_{i(-\infty)}^{i(t)}
\]
在$t=-∞$时，电感电流为零，所以电感吸收的磁场能量为
\[
W_m=\frac{1}{2}Li^2=\frac{\Psi^2}{2L}
\]
式是电感储存磁场能量的计算公式。当电感电流的绝对值增大时，电感吸收能量；电流减小时，释放能量，所以电感元件也是一个储能元件（无源元件）。

\section{独立电源}
电源模型有电压源和电流源两种，均为独立电源（源电压/电流不受电路其他量控制）。

\subsection{电压源}
理想电压源的端口电压与电流大小无关，或为常量或随时间按确定规律变化。

在电路模型中，电压源的符号为“+”、“-”极表示，其端口特性是：无论流过的电流是多少，端口电压恒等于源电压$u_s$。

电压源能够提供确定的电源电压$u_s$，称为源电压。所谓“确定”是指源电压$u_s$与流过电压源的电流无关，电压源的电流将由与其相连的电路来确定。$u_s$可以是常量，称为直流电压源（或恒定电压源），记作$u_s=U_s$；也可以是时变量，记作$u_s=u_s(t)$（如正弦电压源）。源电压$u_s$的方向是从“+”极到“-”极。

电压源的端口特性可以用解析式表示
\[
u=u_s
\]
式中不含电流$i$，表示源电压不受电流的制约。源电压为常量时，端口电压与电流之间的关系可用平行于电流坐标轴的直线表示，称为电源的端口特性曲线。当源电压$u_s=0$即电压源置零时，电压源的作用相当于短路。

在关联参考方向下（电流从“+”极流入），电压源吸收的功率是$p=u_si$；若电流从“+”极流出（非关联参考方向），则$p=-u_si$为发出的功率。当$p>0$时，电压源吸收能量（充电）；$p<0$时，发出能量（供电）。

\subsection{电流源}
理想电流源的端口电流与端口电压无关，恒等于源电流$i_s$。

在电路模型中，电流源的符号为箭头表示电流方向，其端口特性是：无论端口电压是多少，端口电流恒等于源电流$i_s$。

电流源能够提供确定的端口电流$i_s$，称为源电流。这里“确定”是指$i_s$与电流源端口电压无关，电流源的端口电压决定于它所接的外电路。$i_s$可以是常量即恒定电流，记作$i_s=I_s$，也可以是时变量，记作$i_s=i_s(t)$。

电流源端口特性的解析表达式为
\[
i=i_s
\]
式不含电压$u$，表明电流源的端口电流不受端口电压的制约。如果源电流$i_s$为常量，则式对应$u-i$平面上一条平行于电压坐标轴的直线。当$i_s=0$即电流源置零时，电流源的作用相当于开路。

在关联参考方向下（电流从“+”极流入），电流源吸收的功率为$p=ui_s$；非关联参考方向下，发出的功率为$p=-ui_s$。当特性曲线位于第$\mathrm{I}$、第$\mathrm{III}$象限时，$p>0$，电流源处于充电状态；当特性曲线位于第$\mathrm{II}$、第$\mathrm{IV}$象限时，$p<0$，电流源处于供电状态。

电压源和电流源作为元件模型，是实际电源（如电池、发电机）的理想化近似：实际电压源存在内阻，端口电压随电流增大而降低；实际电流源存在内电导，端口电流随电压增大而减小。

\section{受控电源}
在电路模型中还要用到另一类电源，它们的源电压或源电流受电路中另一处的电压或电流控制，这类电源称为受控电源或非独立电源。

受控源是一种用来表示晶体管、运算放大器等电子器件的电路模型，本书只讨论线性受控源，即源电压（流）与控制电压（流）成正比关系。用菱形符号表示受控源，以便与独立源相区分。

线性受控源分为四种类型，其端口方程分别为：
\begin{enumerate}
    \item 电压控制电压源（\textbf{VCVS}）：$\begin{cases}u_s = \alpha u_c \\i_c = 0\end{cases}$（$\alpha$为电压放大系数，无量纲）
    \item 电流控制电压源（\textbf{CCVS}）：$\begin{cases}u_s = r i_c \\u_c = 0\end{cases}$（$r$为转移电阻，单位$\Omega$）
    \item 电压控制电流源（\textbf{VCCS}）：$\begin{cases}i_s = g u_c \\i_c = 0\end{cases}$（$g$为转移电导，单位$\mathrm{S}$）
    \item 电流控制电流源（\textbf{CCCS}）：$\begin{cases}i_s = \beta i_c \\u_c = 0\end{cases}$（$\beta$为电流放大系数，无量纲）
\end{enumerate}
式中带脚注$s$的变量表示源量，带脚注$c$的变量表示控制量。参量$\alpha\text{、}r\text{、}g\text{、}\beta$都是常量，由器件特性决定。

受控源是双端口元件：一个端口是控制端口（输入端口），另一个是受控端口（输出端口）。控制端口若为开路（$i_c=0$），则是电压控制型；若为短路（$u_c=0$），则是电流控制型。

从电路模型角度，受控源本身不消耗/提供能量（理想情况），其能量由控制电路提供，因此受控源属于无源元件（但由受控源组成的电路可能是有源的）。

\section{基尔霍夫定律}
1847年基尔霍夫建立了电路结构对电流、电压的约束关系，分别称为基尔霍夫电流定律（KCL）和基尔霍夫电压定律（KVL），统称基尔霍夫定律，是分析集中参数电路的基本定律。

\subsection{电路结构基本概念}
电路元件通过端子相互连接便形成具有特定结构的电路。本书将每个二端元件称为一条支路，若干支路的连接点称为节点（两个支路的连接点称为简单节点），由支路构成的闭合路径称为回路，内部不含其他支路的回路称为网孔。

一条支路的电流是指流过该支路的电流，支路的电压是指支路两个端子之间的电压。

\subsection{基尔霍夫电流定律（KCL）}
基尔霍夫电流定律表述为：在集中参数电路中，任一时刻流出（或流入）任一节点的支路电流代数和等于零，即
\[
\sum i_k=0
\]
通常规定：电流参考方向为流出节点时取正号，流入节点时取负号（也可反之，只需统一）。

式称为节点的$\mathrm{KCL}$方程，它反映了与任意一个节点相连的所有支路电流所满足的电荷守恒约束（节点不能储存电荷）。

$\mathrm{KCL}$还可推广用于流出（或流入）闭合边界的支路电流（称为“广义节点”），表述为：在集中参数电路中，任一时刻流出（或流入）任一闭合边界$S$的支路电流代数和等于零，即
\[
\sum_{k\in S} i_k=0
\]
式称为广义基尔霍夫电流定律。

如将流出节点的电流记为正、流入记为负，则$\mathrm{KCL}$也可表述为：任一时刻，流出任一节点（或闭合边界）的电流之和等于流入该节点（或闭合边界）的电流之和，即
\[
\sum i_{\text{流出}}=\sum i_{\text{流入}}
\]

\subsection{基尔霍夫电压定律（KVL）}
基尔霍夫电压定律表述为：在集中参数电路中，任一时刻沿任一回路各支路电压的代数和等于零，即
\[
\sum u_k=0
\]
通常规定：当支路电压的参考方向与回路绕行方向相同时，$u_k$前面取正号；反之取负号。

式称为回路的$\mathrm{KVL}$方程，它是任一回路所包含的支路电压必须满足的能量守恒约束（单位电荷沿回路绕行一周，电场力做功为零）。

推广到一般情况就是：沿任一回路，各支路电压降的代数和等于电压升的代数和，即
\[
\sum u_{\text{降}}=\sum u_{\text{升}}
\]
进一步推广：集中参数电路中，任意两点之间的电压具有确定值，与计算路径无关（电压的单值性）。

对电路中的每个网孔或回路均可列出$\mathrm{KVL}$方程，但并非所有方程都是独立的。可以证明：对于具有$b$条支路、$n$个节点的电路，独立$\mathrm{KVL}$方程的个数为$b-(n-1)$（平面电路的网孔数等于独立$\mathrm{KVL}$方程数）。

基尔霍夫定律仅与电路的拓扑结构有关，与元件的性质无关（线性/非线性、时变/时不变均适用）。
% 第2章 线性直流电路（章节子文件，需被主文件调用）
\chapter{线性直流电路}

\section{电阻网络的等效}
电阻的串联

电阻的并联

电阻的星形联结和三角形联结

电阻的星形联结，简称$\mathrm{Y}$形联结或$\mathrm{T}$形联结，。三角形联结，。

可归纳为
\[
R_{1,2,3}=\frac{\boldsymbol{\Delta}\text{形联结中与该节点相连两电阻之积}}{\boldsymbol{\Delta}\text{形联结中各电阻之和}}
\]

可归纳为
\[
R_{1,2,3}=\frac{\mathrm{Y}\text{形联结各电阻两两乘积之和}}{\mathrm{Y}\text{形联结中不相邻电阻}}
\]

\section{含源支路的等效变换}
图表示，它们的电压、电流关系分别为
\[
U=U_s-R_iI
\]
\[
I=I_s-G_iU
\]

当电路参数满足一定条件时，电压源和电阻串联支路与电流源和电导并联支路可以相互等效。为求出等效条件，将式改写成
\[
U=\frac{1}{G_i}I_s-\frac{1}{G_i}I
\]
比较式与式得到两电路的等效变换条件为
\[
U_s=\frac{I_s}{G_i},\quad R_i=\frac{1}{G_i}
\]
图中电压源

受控电压源与电阻的串联也可等效成受控电流源与同一电阻的并联，或反之

\section{支路电流法}

\section{回路电流法}
回路电流

回路电流方程的列写

\section{节点电压法}

\section{运算放大器}

\section{含运算放大器电路的分析}
反相放大器

同相放大器

加法器

差分放大器
% 第3章 电路定律
\chapter{电路定律}

\section{置换定理}
置换定理：在任意线性和非线性电路中，若已知某一端口的电压和电流分别为$U$和$I$，则可用$U_S=U$的电压源或$I_S=I$的电流源来置换此一端口，而不影响电路中其他部分的电压和电流。

在应用置换定理时还需说明：
（1）置换定理要求置换后的电路有唯一解
（2）除被置换部分发生变化外，其余部分在置换前后必须保持完全相同
根据置换定理可以得出推论：

\section{齐性定理和叠加定理}
\subsection{齐性定理}
齐性定理：在只有一个激励$X$作用的线性电路中，设任一响应为$Y$，记作$Y=f(X)$，若将该激励乘以常数$K$，则对应的响应$Y'$也等于原来响应乘以同一常数，即$Y'=f(KX)=Kf(X)=KY$

从数学上看，

对于线性直流电路，其电路方程为线性代数方程，此时齐性定理可直观表述为：若电路中只有一个激励，则响应与激励成正比。

\subsection{叠加定理}
叠加定理：在线性唯一解电路中，有几个独立电源共同作用产生的响应等于各个独立电源单独作用时产生响应的代数叠加。

值得注意的是，功率与独立电源的源电压和源电流不是线性关系，所以功率不能直接用叠加关系来计算。

对于线性直流电路

\section{等效电源定理}
\subsection{戴维宁定理}
线性含源一端口网络的对外作用可以用一个电压源串联电阻的电路来等效代替。其中电源的源电压等于此一端口网络的开路电压，而电阻等于此一端口网络内部独立电源置零后所得无独立源一端口网络的等效电阻。

\subsection{诺顿定理}

\section{特勒根定理}
特勒根定理：在电路

\section{互易定理}
定理（第一种形式）：对于由一个独立电压源和若干线性二端电阻组成的电路，当此电压源

定理（第二种形式）：

定理（第三种形式）：

\section{对偶原理}
如果
% 第4章 正弦电流电路
\chapter{正弦电流电路}

\section{正弦电流}
随正弦规律变化的电流称为正弦电流或正弦交流

\section{正弦量的相量表示法}
复数

正弦量的相量表示
\[
\dot{A}_m = A_m e^{j\psi} = A_m \angle \psi
\]

正弦量运算与相量运算的对应关系
1.唯一性
2.线性性质
3.微分规则

\section{基尔霍夫的相量形式}

\section{RLC元件上电压与电流的关系}
电阻元件

电感元件
可得电感上电压、电流相量关系
\[
\]

电容元件

\section{阻抗与导纳}
阻抗
即阻抗端口电压、电流有效值（或振幅）之比等于阻抗模；端口电压超前于电流的相位差等于阻抗角$\psi$。

导纳

\section{正弦电流电路的相量分析法}

\section{正弦电流电路的功率}
瞬时功率

平均功率

无功功率
定义无功功率
\[
\]

视在功率

功率因数的提高

复功率

\section{最大功率传输定理}

\section{耦合电感}
互感元件端口特性方程

互感元件的等效电路和去耦合等效

\section{含互感元件的正弦电流电路}
含互感电路的等效化简

\section{理想变压器}
% 第5章 三相电路
\chapter{三相电路}

\section{三相电源和三相电路}
如果三相电压不仅频率相同、波形和幅值相同，而且变化进程的时间差相等，即彼此相差，则称为对称三相电压

如果三相电压，其瞬时表达式为
\[
\]
取$k=1$，则为
\[
\]
其波形如图所示，对称正弦三相电压有效值相量为
\[
\]
式中
\[
\]
是一个旋转因子

三相电压的相位。式中三相电压的相位的次序为，称为正序或顺序。

由三相电源、三相负载和连接导线所组成的电路称为三相电路

\section{星形联结和三角形联结}
星形联结
图中三相电源和负载均为星形，或称$\mathrm{Y}$形。三个相连在一起的公共点称为中性点，。中性点之间的连线称为中性线。在这种连接方式中，电源

在对称三相电路中，各相电流必然是对称的。此时
\[
\]
上式说明

端线电流

根据星形联结的电路结构可知：后者用相量表示为
\[
\]

在相电压对称的情况下，线电压与相电压的关系可进一步简化。

同理，可以得出。
\[
\]
图画出了所示电源的对称相电压和线电压的相量图。则有
\[
\]
按照

在对称情况下，三个相电流有效值相等

三角形联结
图中，则三相总电压为
\[
\]
与上式对应的电压相量图。假如，则三相总电压为
\[
\]
相应的相量图。总电压有效值为相电压有效值的二倍，

当把负载接成三角形时，如果

在三角形联结中

在三角形联结中，用相量表示为
\[
\]

在负载端

如果三个相电流是对称的，并由此得线电流与相电流的关系
\[
\]
可见在三角形联结中，当正弦相电流对称时，线电流也是对称的，其有效值等于，即
\[
\]
当按

三角形联结中，在对称情况下可以记作
\[
\]

上面

在三相电路

一个三相负载应当接成

\section{对称三相电路的计算}
图

设，对负载中性点列电压方程
\[
\]
因为回路电压方程为
\[
\]
上式表明

对于比较复杂

\section{不对称三相电路示例}

\section{三相电路的功率}
% 第6章 非正弦周期电流电路
\chapter{非正弦周期电流电路}

\section{非正弦周期电流和电压}
首先

其次

再次

最后

\section{非正弦周期函数的傅里叶级数}
工程上遇到的周期函数如下的傅里叶级数
\[
\]
式中，平均值：
\[
\]
为了求得$a_k$，即
\[
\]
根据数学上，得到
\[
\]
即$a_k$

用类似方法可得
\[
\]

几种常见周期函数的傅里叶级数

\section{非正弦周期量的有效值、平均功率}
有效值
周期量的有效值等于其瞬时值的方均根值，即
任意周期量的有效值等于它的恒定分量

平均功率

\section{非正弦周期电流电路的计算}
% 第7章 频率特性和谐振现象
\chapter{频率特性和谐振现象}

\section{网络函数和频率特性}
若将所关心的网络某一电压或电流作为响应，激励与响应的比例系数记作
\[
H(j\omega) = \frac{\dot{Y}(j\omega)}{\dot{X}(j\omega)}
\]
该系数称为网络函数，其中$\dot{X}(j\omega)$为激励相量，$\dot{Y}(j\omega)$为响应相量。

如果激励和响应均为集中参数正弦量，其网络函数为
\[
H(j\omega) = |H(j\omega)|e^{j\varphi(\omega)}
\]
式中$|H(j\omega)|$为网络函数的模（幅频特性），$\varphi(\omega)$为网络函数的辐角（相频特性）。

可见网络函数的模和辐角均为角频率$\omega$的函数，即
\[
|H(j\omega)| = f(\omega), \quad \varphi(\omega) = \theta(\omega)
\]
由此可得频率特性的完整表达式：
\[
H(j\omega) = |H(j\omega)|\angle\varphi(\omega)
\]
式中幅频特性反映响应与激励的幅值比随频率的变化规律，相频特性反映响应与激励的相位差随频率的变化规律。

\section{串联谐振电路}
根据谐振的定义，RLC串联电路发生谐振时，端口电压与电流同相位，即阻抗的虚部为零：
\[
\mathrm{Im}[Z(j\omega)] = 0
\]
即
\[
\omega L - \frac{1}{\omega C} = 0
\]

由上式解得谐振角频率$\omega_0$，电感电抗$X_L$和电容电抗$X_C$分别为
\[
X_L = \omega_0 L, \quad X_C = \frac{1}{\omega_0 C}
\]
式中$\omega_0$称为$\mathrm{RLC}$串联电路的谐振角频率，其计算公式为：
\[
\omega_0 = \frac{1}{\sqrt{LC}}
\]

发生串联谐振时，令
\[
\rho = \omega_0 L = \frac{1}{\omega_0 C}
\]
$\rho$称为$\mathrm{RLC}$串联电路的特性阻抗。将谐振角频率代入特性阻抗公式，得
\[
\rho = \sqrt{\frac{L}{C}}
\]
又令
\[
Q = \frac{\rho}{R} = \frac{\omega_0 L}{R} = \frac{1}{\omega_0 CR}
\]
$Q$称为$\mathrm{RLC}$串联电路的品质因数。可得品质因数的另外计算公式：
\[
Q = \frac{U_L}{U} = \frac{U_C}{U}
\]
当外加电压$U$不变时，电路中电流$I$、电感电压$U_L$、电容电压$U_C$与角频率$\omega$的关系是
\[
I(\omega) = \frac{U}{\sqrt{R^2 + \left(\omega L - \frac{1}{\omega C}\right)^2}}, \quad U_L = I\omega L, \quad U_C = \frac{I}{\omega C}
\]

根据幅频特性曲线，当电流下降到最大值的$\frac{1}{\sqrt{2}}$倍时，对应的角频率称为截止角频率，记为
\[
\omega_{1,2} = \omega_0 \sqrt{1 \pm \left(\frac{1}{Q}\right)^2}
\]
此时电流满足：
\[
I(\omega_{1,2}) = \frac{I_0}{\sqrt{2}}
\]
\[
I_0 = \frac{U}{R} \quad (\text{谐振时的最大电流})
\]

引用$\mathrm{RLC}$串联电路各物理量的有效值关系：
\[
U = IR, \quad U_L = I\omega L, \quad U_C = \frac{I}{\omega C}
\]
可见谐振时$U_L = U_C = QU$，电容和电感电压为电源电压的$Q$倍，因此串联谐振也称为电压谐振。

如果$\mathrm{RLC}$串联电路中电阻$R$趋于零，则品质因数$Q \to \infty$，谐振时电容和电感电压将远大于电源电压，实际电路中需避免这种情况。

\section{RLC串联电路的频率特性}
当以电阻电压为响应时，转移电压比（网络函数）为
\[
H_R(j\omega) = \frac{\dot{U}_R}{\dot{U}} = \frac{R}{R + j\left(\omega L - \frac{1}{\omega C}\right)}
\]
将谐振角频率$\omega_0 = \frac{1}{\sqrt{LC}}$代入，得
\[
H_R(j\omega) = \frac{1}{1 + jQ\left(\frac{\omega}{\omega_0} - \frac{\omega_0}{\omega}\right)}
\]
\[
|H_R(j\omega)| = \frac{1}{\sqrt{1 + Q^2\left(\frac{\omega}{\omega_0} - \frac{\omega_0}{\omega}\right)^2}}, \quad \varphi(\omega) = -\arctan\left[Q\left(\frac{\omega}{\omega_0} - \frac{\omega_0}{\omega}\right)\right]
\]

由幅频特性可知，通带宽度为截止角频率之差，需满足
\[
|H_R(j\omega)| = \frac{1}{\sqrt{2}}
\]
由此解得下截止角频率和上截止角频率：
\[
\omega_1 = \omega_0 \left(-\frac{1}{2Q} + \sqrt{\left(\frac{1}{2Q}\right)^2 + 1}\right)
\]
\[
\omega_2 = \omega_0 \left(\frac{1}{2Q} + \sqrt{\left(\frac{1}{2Q}\right)^2 + 1}\right)
\]
截止频率与谐振频率、品质因数的关系是
\[
\omega_2 - \omega_1 = \frac{\omega_0}{Q}, \quad \omega_1 \omega_2 = \omega_0^2
\]
可见，$\mathrm{RLC}$带通电路的通带宽度与品质因数成反比，$Q$越大，通带越窄，选频特性越好。

如果以电感电压为响应，其转移电压比为
\[
H_L(j\omega) = \frac{\dot{U}_L}{\dot{U}} = \frac{j\omega L}{R + j\left(\omega L - \frac{1}{\omega C}\right)}
\]
\[
|H_L(j\omega)| = \frac{\omega L}{\sqrt{R^2 + \left(\omega L - \frac{1}{\omega C}\right)^2}}, \quad \varphi_L(\omega) = 90^\circ - \arctan\left(\frac{\omega L - \frac{1}{\omega C}}{R}\right)
\]
幅频特性和相频特性分别为
\[
|H_L(j\omega)| = \frac{Q \cdot \frac{\omega}{\omega_0}}{\sqrt{1 + Q^2\left(\frac{\omega}{\omega_0} - \frac{\omega_0}{\omega}\right)^2}}
\]
\[
\varphi_L(\omega) = 90^\circ - \arctan\left[Q\left(\frac{\omega}{\omega_0} - \frac{\omega_0}{\omega}\right)\right]
\]
当$\omega = \omega_0$时，$|H_L(j\omega_0)| = Q$，$\varphi_L(\omega_0) = 0^\circ$，电感电压与电源电压同相位。

从幅频特性曲线可知，电感电压的最大值出现在$\omega > \omega_0$处，且$Q$越大，最大值越接近$\omega_0$。

（注：此处可插入RLC串联电路电感电压幅频特性曲线图）

最后，以电容电压为响应，其转移电压比为
\[
H_C(j\omega) = \frac{\dot{U}_C}{\dot{U}} = \frac{-j\frac{1}{\omega C}}{R + j\left(\omega L - \frac{1}{\omega C}\right)}
\]
\[
|H_C(j\omega)| = \frac{\frac{1}{\omega C}}{\sqrt{R^2 + \left(\omega L - \frac{1}{\omega C}\right)^2}}, \quad \varphi_C(\omega) = -90^\circ - \arctan\left(\frac{\omega L - \frac{1}{\omega C}}{R}\right)
\]
幅频特性和相频特性分别为
\[
|H_C(j\omega)| = \frac{Q \cdot \frac{\omega_0}{\omega}}{\sqrt{1 + Q^2\left(\frac{\omega}{\omega_0} - \frac{\omega_0}{\omega}\right)^2}}
\]
\[
\varphi_C(\omega) = -90^\circ - \arctan\left[Q\left(\frac{\omega}{\omega_0} - \frac{\omega_0}{\omega}\right)\right]
\]
当$\omega = \omega_0$时，$|H_C(j\omega_0)| = Q$，$\varphi_C(\omega_0) = -180^\circ$，电容电压与电源电压反相位。

（注：此处可插入RLC串联电路电容电压幅频特性曲线图）

\section{并联谐振电路}
由电流源激励的$\mathrm{GCL}$并联电路，其输入导纳为
\[
Y(j\omega) = G + j\left(\omega C - \frac{1}{\omega L}\right)
\]
实现谐振的条件是输入导纳的虚部为零，即端口电压与电流同相位：
\[
\mathrm{Im}[Y(j\omega)] = 0 \implies \omega C - \frac{1}{\omega L} = 0
\]
解得并联谐振角频率：
\[
\omega_0 = \frac{1}{\sqrt{LC}}
\]
与RLC串联电路的谐振角频率相同。
% 第8章 线性动态电路暂态过程的时域分析
\chapter{线性动态电路暂态过程的时域分析}

\section{动态电路的暂态过程}
电容元件

在前面

动态

所以
\[
\]
式中    $i$

消去，得到
\[
\]
式是

一阶

综上所述

电路

\section{电路量的初始值}
1.电容电压$u_c$和电感电流$i_L$初始值的确定
设在线性电容，则有
\[
\]
电容电荷的初始值可以表示为
\[
\]
式中等号右端第一项积分表示$t=0_-$时的电荷$q(0_-)$，故
\[
\]
若在，于是得到
\[
\]
式中，称为电容电荷和电压的原始值。

综上所述，若换路瞬间（）电容电流，在$t=0$时是连续变化的或称渐变。

应用对偶原理，的关系
\[
\]
式成立的条件是，

式及其成立的条件所表示的规律称为换路定律。

上述，即
\[
\]
电路换路将引起电容和电感，吸收或释放的功率为
\[
\]
在

2.除电流初始值的确定
瞬间有
 KCL
 KVL
 电阻元件
 电容元件
 电感元件
式中

\section{一阶电路的零输入响应}
可用一阶

在电路的稳态过程中，

1.$\mathrm{RC}$电路的零输入响应
图所示电路

根据微分方程
\[
\]
根据换路定律得初始条件    $u_c(0_+)=u_c(0_-)=U_0$

方程。其求解步骤为

列出特征方程

特征根

微分方程的通解

式中来确定
\[
\]
将上式的解为
\[
\]
再由电阻
\[
\]
由于

将会产生很大的放电电流。

由。现令
\[
\]
称为$\mathrm{RC}$电路的时间常数

下面。

2.$\mathrm{RL}$电路的零输入响应
图所示电路开关原是断开的。其$\mathrm{KVL}$方程为
\[
\]
根据换路定律，得
\[
\]
式，其特征方程和特征根分别为
\[
\]
\[
\]
微分方程的通解是
\[
\]
式中$=\frac{L}{R}$

称为$\mathrm{RL}$电路的时间常数

积分常数$A$
\[
\]
将，得
\[
\]
根据电阻或电感的元件方程求得
\[
\]
式表明

暂态过程

\section{阶跃函数和冲激函数}
1.单位阶跃函数
考察

单位阶跃函数定义为
\[
\]
无量纲

如果，即
\[
\]
的波形

用，记为
\[
\]
$G(t)$的波形

$G(t)$的幅值为。例如
\[
\]
\[
\]
分别表示

2.单位冲激函数
冲激函数，其定义为
\[
\]
$t=0$时，

单位冲激函数

例如

短路

在图

冲激函数

\section{一阶电路的阶跃响应}
若电路

设方程为
\[
\]
将,得
\[
\]
$u_c$的初始值
\[
\]
方程，即
\[
\]
（1）

图，即
\[
\]
（2）

与方程为
\[
\]
方程，其通解为
\[
\]

（3）

将，得
\[
\]
\[
\]
将，阶跃响应
\[
\]
引用，即
\[
\]
$u_c$

根据，求得电容电流为
\[
\]
根据

$i_c$

线性电路，阶跃特性为
\[
\]
式中

现将响应为
\[
\]
再将，得到
\[
\]
$u_c'$所示

上述，表达式为
\[
\]
当

\section{一阶电路的冲激响应}

\section{正弦电源作用下的一阶电路}

\section{一阶电路的全响应}

\section{一阶电路暂态响应的一般形式}

\section{求解暂态响应的卷积积分法}

\section{二阶电路的暂态过程}

\section{状态变量分析}
% 第9章 线性动态电路暂态过程的复频域分析
\chapter{线性动态电路暂态过程的复频域分析}

\section{拉普拉斯变换}
常见函数的拉普拉斯变换：
\begin{itemize}
    \item 单位阶跃函数：$\mathcal{L}[\varepsilon(t)] = \frac{1}{s} \quad (\text{Re}[s] > 0)$
    \item 指数函数：$\mathcal{L}[e^{at}\varepsilon(t)] = \frac{1}{s-a} \quad (\text{Re}[s] > a)$
    \item 单位冲激函数：$\mathcal{L}[\delta(t)] = 1 \quad (\text{Re}[s] > -\infty)$
\end{itemize}

\section{拉普拉斯变换的基本性质}
\begin{itemize}
    \item 线性性质：$\mathcal{L}[a f_1(t) + b f_2(t)] = a F_1(s) + b F_2(s)$
    \item 微分性质：$\mathcal{L}[f'(t)] = sF(s) - f(0_-)$
    \item 积分性质：$\mathcal{L}\left[\int_{0_-}^t f(\tau)d\tau\right] = \frac{F(s)}{s}$
    \item 延迟性质：$\mathcal{L}[f(t-t_0)\varepsilon(t-t_0)] = e^{-st_0}F(s) \quad (t_0 > 0)$
    \item 位移性质：$\mathcal{L}[e^{-at}f(t)] = F(s+a) \quad (\text{Re}[s+a] > \sigma_0)$
    \item 初值定理：$f(0_+) = \lim\limits_{s \to \infty} sF(s) \quad (\text{满足初值定理条件})$
    \item 终值定理：$f(\infty) = \lim\limits_{s \to 0} sF(s) \quad (\text{满足终值定理条件})$
    \item 卷积定理：$\mathcal{L}[f_1(t) * f_2(t)] = F_1(s)F_2(s)$
\end{itemize}

\section{拉普拉斯逆变换}
拉普拉斯逆变换的求解核心是有理分式分解，按分子分母阶数分两种情况：
\begin{itemize}
    \item 分子阶数 $n > m$（分母阶数）：先做多项式除法，将有理分式分解为多项式 + 真分式
    \item 分子阶数 $n \leq m$（分母阶数）：直接对真分式进行部分分式分解（单极点、共轭复极点、重极点）
\end{itemize}

\section{复频域中的电路定律与电路模型}
\begin{itemize}
    \item 复频域中的基尔霍夫定律：
          \begin{itemize}
              \item KCL：$\sum I_k(s) = 0$（复频域节点电流代数和为零）
              \item KVL：$\sum U_k(s) = 0$（复频域回路电压代数和为零）
          \end{itemize}
    \item 复频域中的元件电压与电流关系：
          \begin{itemize}
              \item 电阻：$U_R(s) = R I_R(s)$
              \item 电感：$U_L(s) = sL I_L(s) - Li_L(0_-)$
              \item 电容：$U_C(s) = \frac{1}{sC} I_C(s) + \frac{u_C(0_-)}{s}$
          \end{itemize}
    \item 复频域电路模型、运算阻抗与运算导纳：
          \begin{itemize}
              \item 运算阻抗：$Z(s) = \frac{U(s)}{I(s)}$（关联参考方向）
              \item 运算导纳：$Y(s) = \frac{1}{Z(s)} = \frac{I(s)}{U(s)}$
          \end{itemize}
\end{itemize}

\section{用拉普拉斯变换分析线性动态电路的暂态过程}
分析步骤：
\begin{enumerate}
    \item 求电路的初始状态 $u_C(0_-)$、$i_L(0_-)$；
    \item 将激励源做拉普拉斯变换，得到复频域激励 $E(s)$；
    \item 建立复频域电路模型（代入初始状态，元件用运算阻抗/导纳表示）；
    \item 利用复频域KCL/KVL列方程，求解响应的复频域表达式 $R(s)$；
    \item 对 $R(s)$ 做拉普拉斯逆变换，得到时域响应 $r(t)$。
\end{enumerate}

\section{网络函数}
\begin{itemize}
    \item 网络函数定义：$H(s) = \frac{R(s)}{E(s)}$（零状态下，响应象函数与激励象函数之比）
    \item 网络函数的极点位置与单位冲激特性的关系：
          \begin{itemize}
              \item 极点在左半复平面：冲激响应衰减（稳定电路）；
              \item 极点在虚轴上：冲激响应等幅振荡；
              \item 极点在右半复平面：冲激响应发散（不稳定电路）。
          \end{itemize}
    \item 复频域网络函数与复数网络函数的关系：$H(j\omega) = H(s)\big|_{s=j\omega}$（$s$ 沿虚轴取值）
\end{itemize}
% 第10章 二端口网络（章节子文件，需被主文件调用）
\chapter{二端口网络}

\section{导纳参数方程和阻抗参数方程}

\subsection{导纳参数方程（Y参数）}
二端口网络的导纳参数（Y参数）以端口电压 $\dot{U}_1\text{、}\dot{U}_2$ 为自变量，端口电流 $\dot{I}_1\text{、}\dot{I}_2$ 为因变量，其代数方程为：
\[
\begin{cases} 
\dot{I}_1 = Y_{11}\dot{U}_1 + Y_{12}\dot{U}_2 \\ 
\dot{I}_2 = Y_{21}\dot{U}_1 + Y_{22}\dot{U}_2 
\end{cases}
\]

写成矩阵形式（Y参数矩阵方程）：
\[
\begin{bmatrix} 
\dot{I}_1 \\ 
\dot{I}_2 
\end{bmatrix} = 
\mathbf{Y}
\begin{bmatrix} 
\dot{U}_1 \\ 
\dot{U}_2 
\end{bmatrix}
\]
其中 $\mathbf{Y} = \begin{bmatrix} Y_{11} & Y_{12} \\ Y_{21} & Y_{22} \end{bmatrix}$ 称为二端口网络的导纳矩阵。

Y参数的物理意义及计算方法（短路参数，下标为0表示短路）：
\[
\begin{aligned}
Y_{11} &= \left. \frac{\dot{I}_1}{\dot{U}_1} \right|_{\dot{U}_2 = 0} \quad (\text{端口2短路时，端口1的输入导纳}) \\
Y_{12} &= \left. \frac{\dot{I}_1}{\dot{U}_2} \right|_{\dot{U}_1 = 0} \quad (\text{端口1短路时，端口2对端口1的转移导纳}) \\
Y_{21} &= \left. \frac{\dot{I}_2}{\dot{U}_1} \right|_{\dot{U}_2 = 0} \quad (\text{端口2短路时，端口1对端口2的转移导纳}) \\
Y_{22} &= \left. \frac{\dot{I}_2}{\dot{U}_2} \right|_{\dot{U}_1 = 0} \quad (\text{端口1短路时，端口2的输入导纳})
\end{aligned}
\]

互易与对称条件：
\begin{itemize}
    \item 互易网络（无受控源）：$Y_{12} = Y_{21}$（转移导纳相等）
    \item 对称网络（电气结构对称）：$Y_{12} = Y_{21}$ 且 $Y_{11} = Y_{22}$（输入导纳相等）
\end{itemize}

\subsection{阻抗参数方程（Z参数）}
二端口网络的阻抗参数（Z参数）以端口电流 $\dot{I}_1\text{、}\dot{I}_2$ 为自变量，端口电压 $\dot{U}_1\text{、}\dot{U}_2$ 为因变量，其代数方程为：
\[
\begin{cases} 
\dot{U}_1 = Z_{11}\dot{I}_1 + Z_{12}\dot{I}_2 \\ 
\dot{U}_2 = Z_{21}\dot{I}_1 + Z_{22}\dot{I}_2 
\end{cases}
\]

写成矩阵形式（Z参数矩阵方程）：
\[
\begin{bmatrix} 
\dot{U}_1 \\ 
\dot{U}_2 
\end{bmatrix} = 
\mathbf{Z}
\begin{bmatrix} 
\dot{I}_1 \\ 
\dot{I}_2 
\end{bmatrix}
\]
其中 $\mathbf{Z} = \begin{bmatrix} Z_{11} & Z_{12} \\ Z_{21} & Z_{22} \end{bmatrix}$ 称为二端口网络的阻抗矩阵。

Z参数的物理意义及计算方法（开路参数，下标为0表示开路）：
\[
\begin{aligned}
Z_{11} &= \left. \frac{\dot{U}_1}{\dot{I}_1} \right|_{\dot{I}_2 = 0} \quad (\text{端口2开路时，端口1的输入阻抗}) \\
Z_{12} &= \left. \frac{\dot{U}_1}{\dot{I}_2} \right|_{\dot{I}_1 = 0} \quad (\text{端口1开路时，端口2对端口1的转移阻抗}) \\
Z_{21} &= \left. \frac{\dot{U}_2}{\dot{I}_1} \right|_{\dot{I}_2 = 0} \quad (\text{端口2开路时，端口1对端口2的转移阻抗}) \\
Z_{22} &= \left. \frac{\dot{U}_2}{\dot{I}_2} \right|_{\dot{I}_1 = 0} \quad (\text{端口1开路时，端口2的输入阻抗})
\end{aligned}
\]

互易与对称条件：
\begin{itemize}
    \item 互易网络（无受控源）：$Z_{12} = Z_{21}$（转移阻抗相等）
    \item 对称网络（电气结构对称）：$Z_{12} = Z_{21}$ 且 $Z_{11} = Z_{22}$（输入阻抗相等）
\end{itemize}

\section{传输参数方程和混合参数方程}

\subsection{传输参数方程（A参数）}
传输参数（A参数，又称T参数）用于描述二端口网络的级联特性，以输出端口的 $\dot{U}_2\text{、}-\dot{I}_2$ 为自变量，输入端口的 $\dot{U}_1\text{、}\dot{I}_1$ 为因变量，其代数方程为：
\[
\begin{cases} 
\dot{U}_1 = A_{11}\dot{U}_2 + A_{12}(-\dot{I}_2) \\ 
\dot{I}_1 = A_{21}\dot{U}_2 + A_{22}(-\dot{I}_2) 
\end{cases}
\]

写成矩阵形式（A参数矩阵方程）：
\[
\begin{bmatrix} 
\dot{U}_1 \\ 
\dot{I}_1 
\end{bmatrix} = 
\mathbf{A}
\begin{bmatrix} 
\dot{U}_2 \\ 
-\dot{I}_2 
\end{bmatrix}
\]
其中 $\mathbf{A} = \begin{bmatrix} A_{11} & A_{12} \\ A_{21} & A_{22} \end{bmatrix}$ 称为传输矩阵。

A参数的物理意义及计算方法：
\[
\begin{aligned}
A_{11} &= \left. \frac{\dot{U}_1}{\dot{U}_2} \right|_{\dot{I}_2 = 0} \quad (\text{端口2开路时，电压传输比}) \\
A_{21} &= \left. \frac{\dot{I}_1}{\dot{U}_2} \right|_{\dot{I}_2 = 0} \quad (\text{端口2开路时，转移导纳}) \\
A_{12} &= \left. \frac{\dot{U}_1}{-\dot{I}_2} \right|_{\dot{U}_2 = 0} \quad (\text{端口2短路时，转移阻抗}) \\
A_{22} &= \left. \frac{\dot{I}_1}{-\dot{I}_2} \right|_{\dot{U}_2 = 0} \quad (\text{端口2短路时，电流传输比})
\end{aligned}
\]

互易与对称条件：
\begin{itemize}
    \item 互易网络（无受控源）：$\det(\mathbf{A}) = A_{11}A_{22} - A_{12}A_{21} = 1$（传输矩阵的行列式为1）
    \item 对称网络（电气结构对称）：$\det(\mathbf{A}) = 1$ 且 $A_{11} = A_{22}$（电压/电流传输比相等）
\end{itemize}

Z参数转换为A参数的公式（仅适用于互易网络，$Z_{12}=Z_{21}$）：
\[
\begin{aligned}
A_{11} &= \frac{Z_{11}}{Z_{21}}, \quad A_{12} = \frac{Z_{11}Z_{22} - Z_{12}Z_{21}}{Z_{21}}, \\
A_{21} &= \frac{1}{Z_{21}}, \quad A_{22} = \frac{Z_{22}}{Z_{21}}
\end{aligned}
\]

\subsection{混合参数方程（H参数）}
混合参数（H参数）结合了阻抗和导纳的特点，以 $\dot{I}_1\text{、}\dot{U}_2$ 为自变量，$\dot{U}_1\text{、}\dot{I}_2$ 为因变量（适用于晶体管电路分析），其代数方程为：
\[
\begin{cases} 
\dot{U}_1 = H_{11}\dot{I}_1 + H_{12}\dot{U}_2 \\ 
\dot{I}_2 = H_{21}\dot{I}_1 + H_{22}\dot{U}_2 
\end{cases}
\]

写成矩阵形式（H参数矩阵方程）：
\[
\begin{bmatrix} 
\dot{U}_1 \\ 
\dot{I}_2 
\end{bmatrix} = 
\mathbf{H}
\begin{bmatrix} 
\dot{I}_1 \\ 
\dot{U}_2 
\end{bmatrix}
\]
其中 $\mathbf{H} = \begin{bmatrix} H_{11} & H_{12} \\ H_{21} & H_{22} \end{bmatrix}$ 称为混合参数矩阵。

H参数的物理意义及计算方法：
\[
\begin{aligned}
H_{11} &= \left. \frac{\dot{U}_1}{\dot{I}_1} \right|_{\dot{U}_2 = 0} \quad (\text{端口2短路时，端口1的输入阻抗，又称输入电阻}) \\
H_{21} &= \left. \frac{\dot{I}_2}{\dot{I}_1} \right|_{\dot{U}_2 = 0} \quad (\text{端口2短路时，电流放大系数，又称电流传输比}) \\
H_{12} &= \left. \frac{\dot{U}_1}{\dot{U}_2} \right|_{\dot{I}_1 = 0} \quad (\text{端口1开路时，电压反馈系数}) \\
H_{22} &= \left. \frac{\dot{I}_2}{\dot{U}_2} \right|_{\dot{I}_1 = 0} \quad (\text{端口1开路时，端口2的输出导纳})
\end{aligned}
\]

互易与对称条件：
\begin{itemize}
    \item 互易网络（无受控源）：$H_{12} = -H_{21}$（反馈系数与电流放大系数符号相反）
    \item 对称网络（电气结构对称）：$H_{12} = -H_{21}$ 且 $\det(\mathbf{H}) = H_{11}H_{22} - H_{12}H_{21} = 1$
\end{itemize}

\section{二端口网络的等效电路}
根据二端口网络的参数，可构造对应的等效电路（仅适用于互易网络）：
\begin{itemize}
    \item Z参数等效T形电路：$Z_1=Z_{11}-Z_{12},\ Z_2=Z_{22}-Z_{21},\ Z_3=Z_{12}$（串联支路阻抗）
    \item Y参数等效Π形电路：$Y_1=Y_{11}+Y_{12},\ Y_3=-Y_{21}=-Y_{12},\ Y_2=Y_{22}+Y_{21}$（并联支路导纳）
    \item A参数等效T形阻抗电路：$Z_1=\frac{A_{11}-1}{A_{21}},\ Z_2=\frac{A_{22}-1}{A_{21}},\ Z_3=\frac{1}{A_{21}}$
    \item A参数等效Π形导纳电路：$Y_a=\frac{A_{22}-1}{A_{12}},\ Y_b=\frac{A_{11}-1}{A_{12}},\ Y_c=\frac{1}{A_{12}}$
\end{itemize}

\section{二端口网络与电源和负载的连接}
当二端口网络输入端口接电源（内阻 $Z_s$）、输出端口接负载（阻抗 $Z_L$）时，需分析以下参数：
\begin{itemize}
    \item 输入阻抗：$Z_i = \frac{\dot{U}_1}{\dot{I}_1}$（负载 $Z_L$ 接在输出端口时，输入端口的等效阻抗）
    \item 输出阻抗：$Z_o = \frac{\dot{U}_2}{\dot{I}_2}$（电源内阻 $Z_s$ 接在输入端口时，输出端口的等效阻抗）
    \item 特性阻抗（对称二端口）：$Z_c = \sqrt{Z_{11}Z_{22} - Z_{12}^2}$（无反射匹配阻抗）
    \item 对称二端口的传输系数：$\gamma = \frac{1}{2}\ln(A_{11} + A_{12}/Z_c)$（衰减和相移特性）
\end{itemize}

\section{二端口网络的级联}
多个二端口网络级联时，总传输矩阵等于各级联网络传输矩阵的乘积：
\[
\mathbf{A}_{\text{总}} = \mathbf{A}_1 \cdot \mathbf{A}_2 \cdot \dots \cdot \mathbf{A}_n
\]
其中 $\mathbf{A}_1\text{、}\mathbf{A}_2\text{、}\dots\text{、}\mathbf{A}_n$ 为各级联网络的传输矩阵，级联顺序与矩阵相乘顺序一致（从输入到输出）。

级联的优势：无需重新计算总参数，仅通过矩阵乘法即可得到整体特性，适用于滤波器、放大器等多级电路分析。
\chapter{第11章（待补充）}
% 这里可以后续补充第11章的内容
\chapter{非线性电路}

\section{非线性电阻元件的特性}

\section{非线性直流电路方程}

\section{数值分析法}

\section{分段线性近似法}

\section{图解法}

\section{小信号分析法}
\chapter{均匀传输线}

\section{均匀传输线及其方程}
\begin{itemize}
    \item 电路参数的分布性
    \item 均匀线方程：
    \[
    \begin{cases} 
    -\dfrac{\partial u}{\partial x} = R_0 i + L_0 \dfrac{\partial i}{\partial t} \\ 
    -\dfrac{\partial i}{\partial x} = G_0 u + C_0 \dfrac{\partial u}{\partial t} 
    \end{cases}
    \]
\end{itemize}

\section{均匀传输线方程的正弦稳态解}
正弦稳态下相量形式：
\[
\begin{cases} 
-\dfrac{\mathrm{d}\dot{U}(x)}{\mathrm{d}x} = \left(R_0 + \mathrm{j}\omega L_0\right)\dot{I}(x) = Z_0 \dot{I}(x) \\ 
-\dfrac{\mathrm{d}\dot{I}(x)}{\mathrm{d}x} = \left(G_0 + \mathrm{j}\omega C_0\right)\dot{U}(x) = Y_0 \dot{U}(x) 
\end{cases}
\]

二阶微分方程：
\[
\begin{cases} 
\dfrac{\mathrm{d}^2 \dot{U}(x)}{\mathrm{d}x^2} = Z_0 Y_0 \dot{U}(x) = \gamma^2 \dot{U}(x) \\ 
\dfrac{\mathrm{d}^2 \dot{I}(x)}{\mathrm{d}x^2} = Z_0 Y_0 \dot{I}(x) = \gamma^2 \dot{I}(x) 
\end{cases}
\]

传播常数：
\[
\gamma = \sqrt{\left(R_0 + \mathrm{j}\omega L_0\right)\left(G_0 + \mathrm{j}\omega C_0\right)} = \alpha + \mathrm{j}\beta
\]

通解：
\[
\begin{cases} 
\dot{U}(x) = A_1 \mathrm{e}^{-\gamma x} + A_2 \mathrm{e}^{\gamma x} \\ 
\dot{I}(x) = B_1 \mathrm{e}^{-\gamma x} + B_2 \mathrm{e}^{\gamma x} 
\end{cases}
\]

特性阻抗（波阻抗）：
\[
Z_{\mathrm{c}} = \sqrt{\frac{R_0 + \mathrm{j}\omega L_0}{G_0 + \mathrm{j}\omega C_0}}
\]

边界条件（始端$U_1,I_1$）：
\[
\begin{cases} 
\dot{U}(x) = \dot{U}_1 \cosh \gamma x - \dot{I}_1 Z_{\mathrm{c}} \sinh \gamma x \\ 
\dot{I}(x) = -\dfrac{\dot{U}_1}{Z_{\mathrm{c}}} \sinh \gamma x + \dot{I}_1 \cosh \gamma x 
\end{cases}
\]

矩阵形式：
\[
\begin{bmatrix} 
\dot{U}(x) \\ 
\dot{I}(x) 
\end{bmatrix} = 
\begin{bmatrix} 
\cosh \gamma x & -Z_{\text{c}} \sinh \gamma x \\ 
-\dfrac{1}{Z_{\text{c}}} \sinh \gamma x & \cosh \gamma x 
\end{bmatrix}
\begin{bmatrix} 
\dot{U}_1 \\ 
\dot{I}_1 
\end{bmatrix}
\]

\section{均匀传输线上的行波}
\subsection{行波}
电压电流分解为正向/反向行波：
\[
\dot{U}(x) = \dot{U}^+(x) + \dot{U}^-(x),\quad \dot{I}(x) = \dot{I}^+(x) - \dot{I}^-(x)
\]

相速：$v_p=\frac{\omega}{\beta}$，波长：$\lambda=\frac{2\pi}{\beta}$

\subsection{波的反射与终端匹配}
终端反射系数：
\[
N_2 = \frac{Z_L - Z_c}{Z_L + Z_c}
\]

匹配条件：$Z_L=Z_c$（无反射波）

输入阻抗：
\[
Z_i = Z_c \frac{Z_L \cosh \gamma l + Z_c \sinh \gamma l}{Z_L \sinh \gamma l + Z_c \cosh \gamma l}
\]

\section{无损耗线、驻波}
无损耗线（$R_0=0,G_0=0$）：
\begin{itemize}
    \item 传播常数：$\gamma = \mathrm{j}\omega\sqrt{L_0C_0} = \mathrm{j}\beta$
    \item 波阻抗：$Z_c = \sqrt{\frac{L_0}{C_0}}$（纯实数）
    \item 终端开路/短路/接负载的电压电流特性
\end{itemize}

\section{均匀传输线的集中参数等效电路}

\section{无损均匀线方程的通解}

\section{无损均匀线上波的发出}

\section{无损均匀线上波的反射}
\begin{itemize}
    \item 终端开路
    \item 终端短路
    \item 终端接负载
    \item 波的反射规律
\end{itemize}
% 电路理论基础公式汇总 - 修复版
\chapter{电路理论基础公式汇总}
\thispagestyle{plain}  % 恢复页码样式

\section{第一章：电路元件与电路基本定律}

\subsection{电流、电压与电功率}
\begin{enumerate}
    \item 电流定义：$i=\frac{dq}{dt} \formulaNote{（单位：A，库/秒）}$
    \item 电压定义：$u_{ab}=\frac{dW_{ab}}{dq} \formulaNote{（单位：V，焦/库）}$
    \item 电源电动势：$e=\frac{dW}{dq} \formulaNote{（单位：V，非静电力做功）}$
    \item 电功率定义：$p=\frac{dW}{dt}=ui \formulaNote{（关联参考方向，单位：W）}$
    \item 电路能量计算：$W(t)=\int_{t_0}^tp(\xi)d\xi=\int_{t_0}^tu(\xi)i(\xi)d\xi \formulaNote{（单位：J）}$
\end{enumerate}

\subsection{电阻元件}
\begin{itemize}
    \item 欧姆定律：$u=Ri \quad \text{或} \quad i=Gu \formulaNote{（关联参考方向）}$
    \item 电导与电阻关系：$RG=1 \formulaNote{（G为电导，单位：S）}$
    \item 电阻功率：$p=ui=Ri^2=Gu^2 \formulaNote{（耗能元件，$p\ge0$）}$
    \item 电阻总耗能：
    \[
    W=\int_{-\infty}^tp(\xi)d\xi=\int_{-\infty}^tu(\xi)i(\xi)d\xi = R\int_{-\infty}^ti^2(\xi)d\xi=G\int_{-\infty}^tu^2(\xi)d\xi
    \]
\end{itemize}

\subsection{电容元件}
\begin{itemize}
    \item 电容基本公式：$C=\frac{\varepsilon A}{d} \formulaNote{（平行板电容，$\varepsilon$为介电常数）}$
    \item 电荷-电压关系：$q=Cu \formulaNote{（单位：F，库/伏）}$
    \item 电流-电压微分关系：$i=\frac{dq}{dt}=C\frac{du}{dt} \formulaNote{（直流时$i=0$，等效开路）}$
    \item 电荷积分形式：$q(t)=\int_{-\infty}^ti(\xi)d\xi$
    \item 电压积分形式：$u(t)=\frac{1}{C}\int_{-\infty}^ti(\xi)d\xi$
    \item 初始条件下电压：$u(t)=u(t_0)+\frac{1}{C}\int_{t_0}^ti(\xi)d\xi \formulaNote{（$u(t_0)$为初始电压）}$
    \item 电容储能：$W_e(t)=\frac{1}{2}Cu^2=\frac{q^2}{2C} \formulaNote{（单位：J，仅与当前电压有关）}$
\end{itemize}

\subsection{电感元件}
\begin{itemize}
    \item 磁链-电流关系：$\Psi=Li \formulaNote{（L为电感，单位：H，韦/安）}$
    \item 电压-电流微分关系：$u=\frac{d\Psi}{dt}=L\frac{di}{dt} \formulaNote{（直流时$u=0$，等效短路）}$
    \item 电流积分形式：$i(t)=\frac{1}{L}\int_{-\infty}^tu(\xi)d\xi$
    \item 初始条件下电流：$i(t)=i(t_0)+\frac{1}{L}\int_{t_0}^tu(\xi)d\xi \formulaNote{（$i(t_0)$为初始电流）}$
    \item 电感储能：$W_m=\frac{1}{2}Li^2=\frac{\Psi^2}{2L} \formulaNote{（单位：J，仅与当前电流有关）}$
\end{itemize}

\subsection{独立电源}
\begin{itemize}
    \item 电压源端口特性：$u=u_s \formulaNote{（电压与电流无关，置零等效短路）}$
    \item 电流源端口特性：$i=i_s \formulaNote{（电流与电压无关，置零等效开路）}$
\end{itemize}

\subsection{受控电源（线性）}
\begin{itemize}
    \item 电压控制电压源（VCVS）：$\begin{cases}u_s = \alpha u_c \\i_c = 0\end{cases} \formulaNote{（$\alpha$：电压放大系数）}$
    \item 电流控制电压源（CCVS）：$\begin{cases}u_s = r i_c \\u_c = 0\end{cases} \formulaNote{（$r$：转移电阻，$\Omega$）}$
    \item 电压控制电流源（VCCS）：$\begin{cases}i_s = g u_c \\i_c = 0\end{cases} \formulaNote{（$g$：转移电导，S）}$
    \item 电流控制电流源（CCCS）：$\begin{cases}i_s = \beta i_c \\u_c = 0\end{cases} \formulaNote{（$\beta$：电流放大系数）}$
\end{itemize}

\subsection{基尔霍夫定律}
\begin{itemize}
    \item 基尔霍夫电流定律（KCL）：$\sum i_k=0 \formulaNote{（节点电流守恒，流出为正）}$
    \item 广义KCL：$\sum_{k\in S} i_k=0 \formulaNote{（闭合边界电流守恒）}$
    \item 基尔霍夫电压定律（KVL）：$\sum u_k=0 \formulaNote{（回路电压守恒，绕行方向一致为正）}$
    \item KVL推广：$\sum u_{\text{电压降}}=\sum u_{\text{电压升}}$
\end{itemize}

\section{第十章：二端口网络核心公式}

\subsection{导纳参数（Y参数，短路参数）}
\begin{itemize}
    \item 代数方程：
    \[
    \begin{cases} 
    \dot{I}_1 = Y_{11}\dot{U}_1 + Y_{12}\dot{U}_2 \\ 
    \dot{I}_2 = Y_{21}\dot{U}_1 + Y_{22}\dot{U}_2 
    \end{cases}
    \]
    \item 矩阵形式：$\begin{bmatrix} \dot{I}_1 \\ \dot{I}_2 \end{bmatrix} = \begin{bmatrix} Y_{11} & Y_{12} \\ Y_{21} & Y_{22} \end{bmatrix}\begin{bmatrix} \dot{U}_1 \\ \dot{U}_2 \end{bmatrix}$
    \item 参数计算：
    \[
    \begin{aligned}
    Y_{11} &= \left. \frac{\dot{I}_1}{\dot{U}_1} \right|_{\dot{U}_2 = 0}, & Y_{12} &= \left. \frac{\dot{I}_1}{\dot{U}_2} \right|_{\dot{U}_1 = 0} \\
    Y_{21} &= \left. \frac{\dot{I}_2}{\dot{U}_1} \right|_{\dot{U}_2 = 0}, & Y_{22} &= \left. \frac{\dot{I}_2}{\dot{U}_2} \right|_{\dot{U}_1 = 0}
    \end{aligned}
    \]
    \item 互易/对称条件：$Y_{12}=Y_{21}$（互易）；$Y_{12}=Y_{21}$且$Y_{11}=Y_{22}$（对称）
\end{itemize}

\subsection{阻抗参数（Z参数，开路参数）}
\begin{itemize}
    \item 代数方程：
    \[
    \begin{cases} 
    \dot{U}_1 = Z_{11}\dot{I}_1 + Z_{12}\dot{I}_2 \\ 
    \dot{U}_2 = Z_{21}\dot{I}_1 + Z_{22}\dot{I}_2 
    \end{cases}
    \]
    \item 矩阵形式：$\begin{bmatrix} \dot{U}_1 \\ \dot{U}_2 \end{bmatrix} = \begin{bmatrix} Z_{11} & Z_{12} \\ Z_{21} & Z_{22} \end{bmatrix}\begin{bmatrix} \dot{I}_1 \\ \dot{I}_2 \end{bmatrix}$
    \item 参数计算：
    \[
    \begin{aligned}
    Z_{11} &= \left. \frac{\dot{U}_1}{\dot{I}_1} \right|_{\dot{I}_2 = 0}, & Z_{12} &= \left. \frac{\dot{U}_1}{\dot{I}_2} \right|_{\dot{I}_1 = 0} \\
    Z_{21} &= \left. \frac{\dot{U}_2}{\dot{I}_1} \right|_{\dot{I}_2 = 0}, & Z_{22} &= \left. \frac{\dot{U}_2}{\dot{I}_2} \right|_{\dot{I}_1 = 0}
    \end{aligned}
    \]
    \item 互易/对称条件：$Z_{12}=Z_{21}$（互易）；$Z_{12}=Z_{21}$且$Z_{11}=Z_{22}$（对称）
\end{itemize}

\subsection{传输参数（A参数，T参数）}
\begin{itemize}
    \item 代数方程：
    \[
    \begin{cases} 
    \dot{U}_1 = A_{11}\dot{U}_2 + A_{12}(-\dot{I}_2) \\ 
    \dot{I}_1 = A_{21}\dot{U}_2 + A_{22}(-\dot{I}_2) 
    \end{cases}
    \]
    \item 矩阵形式：$\begin{bmatrix} \dot{U}_1 \\ \dot{I}_1 \end{bmatrix} = \begin{bmatrix} A_{11} & A_{12} \\ A_{21} & A_{22} \end{bmatrix}\begin{bmatrix} \dot{U}_2 \\ -\dot{I}_2 \end{bmatrix}$
    \item 互易/对称条件：$A_{11}A_{22}-A_{12}A_{21}=1$（互易）；$A_{11}A_{22}-A_{12}A_{21}=1$且$A_{11}=A_{22}$（对称）
    \item Z转A参数：$A_{11}=\frac{Z_{11}}{Z_{21}},\ A_{12}=\frac{Z_{11}Z_{22}-Z_{12}Z_{21}}{Z_{21}},\ A_{21}=\frac{1}{Z_{21}},\ A_{22}=\frac{Z_{22}}{Z_{21}}$
\end{itemize}

\subsection{混合参数（H参数）}
\begin{itemize}
    \item 代数方程：
    \[
    \begin{cases} 
    \dot{U}_1 = H_{11}\dot{I}_1 + H_{12}\dot{U}_2 \\ 
    \dot{I}_2 = H_{21}\dot{I}_1 + H_{22}\dot{U}_2 
    \end{cases}
    \]
    \item 矩阵形式：$\begin{bmatrix} \dot{U}_1 \\ \dot{I}_2 \end{bmatrix} = \begin{bmatrix} H_{11} & H_{12} \\ H_{21} & H_{22} \end{bmatrix}\begin{bmatrix} \dot{I}_1 \\ \dot{U}_2 \end{bmatrix}$
    \item 互易/对称条件：$H_{12}=-H_{21}$（互易）；$H_{12}=-H_{21}$且$H_{11}H_{22}-H_{12}H_{21}=1$（对称）
\end{itemize}

\subsection{二端口等效电路}
\begin{itemize}
    \item Z参数T形等效：$Z_1=Z_{11}-Z_{12},\ Z_2=Z_{22}-Z_{21},\ Z_3=Z_{12}$
    \item Y参数Π形等效：$Y_1=Y_{11}+Y_{12},\ Y_3=-Y_{21}=-Y_{12},\ Y_2=Y_{22}+Y_{21}$
    \item 级联总传输矩阵：$\mathbf{A}_{\text{总}} = \mathbf{A}_1 \cdot \mathbf{A}_2 \cdot \dots \cdot \mathbf{A}_n$
\end{itemize}

% 补充通用说明，提升实用性
\section{公式使用说明}
\begin{enumerate}
    \item 所有公式默认采用\textbf{关联参考方向}（电压正极流入，电流流出），非关联方向需加负号；
    \item 交流量用$\dot{U}$、$\dot{I}$表示相量，直流量用$U$、$I$表示；
    \item 互易条件仅适用于无受控源的二端口网络，对称条件需电气结构对称；
    \item 电容/电感的积分公式中，$t=-\infty$表示初始状态为零（未充电/无电流）。
\end{enumerate}

% 文档结束
\end{document}